\hypertarget{execution}{%
\chapter{7. Execution}\label{execution}}

Chapter 6 settled what a kennel may read and write. Running a program is
a separate grant. A binary the kennel can read is a file like any other
until it is launched, and whether it may be launched is gated by a
second Landlock access laid over the first, on an allowlist of its own,
and what follows is that allowlist and the honest account of what it
controls and what it does not.

\hypertarget{what-it-gates}{%
\section{7.1 What it gates}\label{what-it-gates}}

Execution control governs which programs a kennel may launch, by gating
\texttt{execve} and its variants on a deny-by-default path allowlist. An
empty allowlist runs nothing, so a bare base template can start no
program at all and a derived template adds exactly the binaries it
needs. The grant is on the act of launching a new program, not on the
file's existence: the binary is present and readable through its
filesystem grant, and the execution layer decides only whether the
kernel will turn it into a running process.

The boundary is narrower than it first sounds, and the design states the
narrow version rather than the flattering one. What is controlled is
\texttt{execve}, the spawning of a new program. It is not a control over
the code a running program then loads. A program the kennel is allowed
to run can map and execute any file it can read, because the kernel
checks the execute access at \texttt{execve} and not at the
\texttt{mmap} a dynamic loader uses to pull in code. The honest claim is
that a kennel cannot spawn a program that is not on its allowlist, not
that it cannot run code that is not on its allowlist. For most workloads
the two are close; for an interpreter they are the same statement read
twice, since an interpreter evaluates code by design. The execution
layer answers which binaries may launch, and the filesystem, network,
and local-service layers answer what a launched binary may then do.

\hypertarget{the-allowlist-and-no_new_privs}{%
\section{7.2 The allowlist and
no\_new\_privs}\label{the-allowlist-and-no_new_privs}}

An execution policy is a shell and an allowlist of paths:

\begin{kennelcode}
[exec]
shell = "/bin/sh"
allow = ["/bin/sh", "/bin/dash", "/usr/bin/dbus-send"]
\end{kennelcode}

The \texttt{allow} list is the execute allowlist, and \texttt{shell}
names the login shell the synthetic passwd gives the masked user. The
allowlist is enforced by Landlock, with
\texttt{LANDLOCK\_ACCESS\_FS\_EXECUTE} on the granted paths and a denial
on everything else, which is why the kernel floor for an execution
policy is 6.10, where that access has its full semantics; on an older
kernel the policy is refused rather than applied at reduced fidelity
(\texttt{refuse-to-start}). The enforcement is independent of
\texttt{\$PATH}: a kennel that invokes a tool by bare name fails the
\texttt{execve} if the resolved binary is not granted, no matter where
it sits on the search path, and the kennel's \texttt{\$PATH} is set only
so the failure reads as a missing tool rather than a mysterious lookup.
The binaries a confinement most wants to keep out, the
privilege-changing ones, are kept out by simply never appearing on any
allowlist: the escalation tools are not denied by name, they are absent
from the grant, so they do not run (\texttt{mediate-use-not-reach}).

Beneath the allowlist sits a second mechanism that does not depend on
it. \texttt{PR\_SET\_NO\_NEW\_PRIVS} is set unconditionally in every
kennel, before the workload runs, and no policy can turn it off. With it
set, a setuid or setgid binary does not gain its owner's identity on
\texttt{execve}, and file capabilities do not apply, so even an
escalation binary that somehow were on the allowlist could not lift the
workload above the operator it runs as (T3.1). This is the
execution-time face of keeping identity without authority: the workload
cannot climb back to privilege through a binary it finds, because the
kernel refuses the climb (\texttt{split-the-uid}). The flag is cheap, it
is non-negotiable, and a framework that let a policy clear it would be
misnamed.

\hypertarget{what-fs_execute-checks}{%
\section{7.3 What FS\_EXECUTE checks}\label{what-fs_execute-checks}}

The execute access is checked at \texttt{execve}, on each file the
kernel opens for execution, and for a dynamically-linked program that is
two files, not one: the binary itself, and the ELF interpreter named in
its \texttt{PT\_INTERP}, the dynamic loader the kernel opens to run it.
An allowlisted dynamic binary therefore needs the execute grant on its
loader as well as on itself, or it does not start. It does not need, and
does not get, an execute grant on the shared libraries the loader then
pulls in. Those arrive through a memory mapping, and Landlock has no
hook on \texttt{mmap} or \texttt{mprotect}, so a file mapped executable
is never checked against the execute access at all. A library loads on
read access alone, and no execute grant on it would change that.

The design follows that fact rather than fighting it. At compile time
the toolchain reads each allowlisted binary's ELF with a parser rather
than by running it, takes the loader path from \texttt{PT\_INTERP}, and
settles the distinct loaders into the signed policy as their own fixed,
auditable set:

\begin{kennelcode}
[exec]
allow   = ["/usr/bin/python3"]
loaders = ["/lib64/ld-linux-x86-64.so.2"]
\end{kennelcode}

The operator wrote the \texttt{allow} line; the \texttt{loaders} set is
the compiler's, \texttt{python3}'s \texttt{PT\_INTERP} resolved and
recorded so the execute check at \texttt{execve} finds the loader
granted as well. A statically-linked binary names no interpreter and
contributes none. The libraries are simply readable, covered by the
ordinary filesystem grants that already reach the system library
directories, and there is deliberately no execute-allowlist over them. A
list of permitted libraries would be a policy the kernel cannot enforce,
since the access it would name is never checked, and a control that
cannot be enforced is worse than no control, because it reads as a
promise the system does not keep.

\hypertarget{writable-and-interpreted}{%
\section{7.4 Writable, and interpreted}\label{writable-and-interpreted}}

Two limits on the allowlist are worth stating because a reader will
otherwise assume them away. The first is that a writable path must not
also be executable. Without that rule, a kennel granted write into a
project and the right to run an interpreter could write a fresh binary
and run it, or write a script and feed it to the interpreter, and the
allowlist would have bought nothing. The \texttt{deny\_writable}
invariant closes it the way a \texttt{noexec} mount would, making the
intersection of writable and executable paths empty by enforcement, so a
binary placed in a writable directory cannot be launched from it.

The second is the interpreter caveat, and it is the honest boundary
again at one remove. Granting an interpreter the right to run grants it
the right to run anything it can read: a kennel whose only executable is
a scripting language can execute any program written in that language,
from a file, from standard input, from the network. The execution grant
gates which interpreters start, not what they do once started, and a
reader who expects an allowlisted interpreter to mean a single permitted
script has misread the control. The containment for what an interpreter
does lives in the other resource classes, not here. Where the workload's
own first binary must be exactly the expected one, a content pin closes
the remaining gap: the policy may carry a set of accepted digests for
the initial binary, and the daemon verifies the resolved binary's hash
against them before the handoff into the kennel, refusing the run on a
mismatch. That check is on the bytes, made by the daemon and never by
the init inside the kennel, and it is the one place the execution
control reaches past the path to the content behind it.

\hypertarget{composing-the-allowlist}{%
\section{7.5 Composing the allowlist}\label{composing-the-allowlist}}

The allowlist in §7.2 was hand-listed, and for a one-off policy that is
the honest form. In practice the same few binaries recur, an interpreter
and its package tools, a version-control client, a build toolchain, and
hand-listing them in every policy invites the drift a confinement exists
to prevent. The recurring grants live in a fragment instead, a signed,
version-pinned bundle a policy includes rather than copies:

\begin{kennelcode}
# fragments/lang-python/policy.toml
name = "lang-python"

[[exec.allow.add]]
path = "/usr/bin/python3"
reason = "Python interpreter"

[[exec.allow.add]]
path = "/usr/bin/pip3"
reason = "Python package installer"
\end{kennelcode}

A fragment is additive: every entry is a \texttt{{[}{[}*.add{]}{]}}
delta appended to the effective policy, so a fragment can only widen a
grant and never remove or override one, which keeps composition
order-independent. The \texttt{lang-python} fragment carries more than
its exec grants, the writable \texttt{\textasciitilde{}/.cache/pip} the
installer needs and the \texttt{pypi.org} proxy entry the index needs,
so one include grants a Python workflow everything it touches across
exec, filesystem, and network at once. A leaf pulls it in by versioned
reference:

\begin{kennelcode}
name = "my-agent"
template_base = "ai-coding-strict@v1"
include = ["lang-python@v1", "vcs-git@v1"]
\end{kennelcode}

Each include is resolved and signature-verified the way the parent
template is, and byte-pinned in the lockfile, so the resolved set
reproduces or the build fails. The loaders of §7.3 do not care where a
binary's grant came from: the compiler resolves \texttt{python3}'s
\texttt{PT\_INTERP} into the settled policy whether the operator
hand-listed it or included \texttt{lang-python}. Two includes that add
conflicting rules for the same destination fail to compile with an
explicit conflict rather than a silent last-wins, and the resolution
machinery that enforces all of this is the compiler's.
